% Section 10.2, from lectures/L10/README.md: the goals, the evaluation questions, and the next
% lecture.

\needspace{10\baselineskip}
\section{Review}
\label{c10:sec:review}

\subsection*{What you should be able to do now}
After this chapter, you should:
\begin{itemize}
  \item Explain how \code{static} restricts a function's visibility to its own file, and why that's
    a different kind of encapsulation than hiding struct fields.
  \item Explain what an opaque struct is and why hiding a struct's definition in a \file{.c} file
    protects its invariants.
  \item Be able to declare, assign, and call through a function pointer, including one stored in a
    struct field.
  \item Be able to design a small vtable-style interface and provide more than one implementation
    of it.
  \item Understand, at a high level, how this manual pattern relates to virtual dispatch in
    higher-level languages.
\end{itemize}

\needspace{6\baselineskip}
\subsection*{Questions to test yourself}
\begin{enumerate}
  \item What stops calling code from calling the \code{static} \code{gpio_write()} in
    \secref{c10:a:2}'s \file{gpio.c} directly by name, and what does it \emph{not} stop it from
    doing to \code{led1.write} or \code{led1.pin}?
  \item Why can't code outside \file{timer.c} read or write a \code{timer_t}'s fields directly,
    once \code{timer_t} is made opaque?
  \item What does \code{void (*on_press)(void);} declare, and how would you call it?
  \item Why does the pin-change ISR only set a flag instead of calling the callback directly?
  \item What problem does a vtable-style interface solve that hardcoding a single GPIO
    implementation doesn't?
\end{enumerate}

\needspace{6\baselineskip}
\subsection*{Looking ahead}
\begin{itemize}
  \item This closes the driver-design arc (Chapters~\ref{c9:ch} and~\ref{c10:ch}), built on the
    peripherals covered in Chapters~\ref{c1:ch}--\ref{c8:ch}.
  \item Everything here, opaque structs, function pointers, vtables, was C standing in for object
    orientation. The next step is Modern Embedded C++, where these patterns become language
    features (classes, access specifiers, virtual functions) instead of conventions enforced by
    hand.
  \item Covered in a separate course, \emph{Modern Embedded C++}:
    \url{https://github.com/qrtech-academy/modern-embedded-cpp}.
\end{itemize}
